\documentclass{ximera}

\begin{document}
	\author{Stitz-Zeager}
	\xmtitle{Answers for InverseFunctions}{}

\mfpicnumber{1} \opengraphsfile{ExercisesforInverseFunctions} % mfpic settings added 

\label{ExercisesforInverseFunctions}




\begin{multicols}{2}
\begin{enumerate}
\addtocounter{enumi}{8}

\item $f^{-1}(x) = \dfrac{x + 2}{6}$
\item $f^{-1}(x) = 42-x$

\setcounter{HW}{\value{enumi}}
\end{enumerate}
\end{multicols}

\begin{multicols}{2}
\begin{enumerate}
\setcounter{enumi}{\value{HW}}

\item  $g^{-1}(t) = 3t-10$
\item $g^{-1}(t)  = -\frac{5}{3} t + \frac{1}{3}$


\setcounter{HW}{\value{enumi}}
\end{enumerate}
\end{multicols}

\begin{multicols}{2}
\begin{enumerate}
\setcounter{enumi}{\value{HW}}


\item $f^{-1}(x) = \frac{1}{3}(x-5)^2+\frac{1}{3}$, $x \geq 5$
\item $f^{-1}(x) = (x - 2)^{2} + 5, \; x \leq 2$

\setcounter{HW}{\value{enumi}}
\end{enumerate}
\end{multicols}

\begin{multicols}{2}
\begin{enumerate}
\setcounter{enumi}{\value{HW}}


\item $g^{-1}(t) = \frac{1}{9}(t+4)^2+1$, $t \geq -4$

\item $g^{-1}(t) = \frac{1}{8}(t-1)^2-\frac{5}{2}$, $t \leq 1$

\setcounter{HW}{\value{enumi}}
\end{enumerate}
\end{multicols}

\begin{multicols}{2}
\begin{enumerate}
\setcounter{enumi}{\value{HW}}

\item $f^{-1}(x) = \frac{1}{3} x^{5} + \frac{1}{3}$
\item $f^{-1}(x) = -(x-3)^3+2$

\setcounter{HW}{\value{enumi}}
\end{enumerate}
\end{multicols}

\begin{multicols}{2}
\begin{enumerate}
\setcounter{enumi}{\value{HW}}

\item $g^{-1}(t) = 5 + \sqrt{t+25}$
\item $g^{-1}(t) = -\sqrt{\frac{t + 5}{3}} - 4$

\setcounter{HW}{\value{enumi}}
\end{enumerate}
\end{multicols}

\begin{multicols}{2}
\begin{enumerate}
\setcounter{enumi}{\value{HW}}

\item $f^{-1}(x) = 3 - \sqrt{x+4}$
\item $f^{-1}(x) =-\frac{\sqrt{x}+1}{2}$, $x > 1$

\setcounter{HW}{\value{enumi}}
\end{enumerate}
\end{multicols}

\begin{multicols}{2}
\begin{enumerate}
\setcounter{enumi}{\value{HW}}

\item $g^{-1}(t) = \dfrac{4t-3}{t}$
\item $g^{-1}(t) = \dfrac{t}{3t+1}$

\setcounter{HW}{\value{enumi}}
\end{enumerate}
\end{multicols}

\begin{multicols}{2}
\begin{enumerate}
\setcounter{enumi}{\value{HW}}

\item $f^{-1}(x) = \dfrac{4x+1}{2-3x}$
\item $f^{-1}(x) = \dfrac{6x + 2}{3x - 4}$

\setcounter{HW}{\value{enumi}}
\end{enumerate}
\end{multicols}

\begin{multicols}{2}
\begin{enumerate}
\setcounter{enumi}{\value{HW}}

\item $g^{-1}(t) = \dfrac{-3t - 2}{t + 3}$
\item $g^{-1}(t) = \dfrac{t-2}{2t-1}$ 

\setcounter{HW}{\value{enumi}}
\end{enumerate}
\end{multicols}

\begin{enumerate}
\setcounter{enumi}{\value{HW}}

\item

\begin{enumerate}

\item  None of the first coordinates of the ordered pairs in $F$ are repeated, so $F$ is a function and none of the second coordinates of the ordered pairs of $F$ are repeated, so $F$ is one-to-one.   $F^{-1} = \{ (0,0), (1,1), (-1,2), (2,3), (-2,4), (3,5), (-3,6)  \}$

\item  Because of the `$\ldots$' it is helpful to determine a formula for the matching. For the even numbers $n$, $n = 0, 2, 4, \ldots$, the ordered pair $\left(n, -\frac{n}{2} \right)$ is in $G$.  For the odd numbers  $n = 1, 3, 5, \ldots$, the ordered pair $\left(n, \frac{n+1}{2} \right)$ is in $G$.  Hence, given any input to $G$, $n$, whether it be even or odd, there is only one output from $G$, either $-\frac{n}{2}$ or $\frac{n+1}{2}$, both of which are functions of $n$. To show $G$ is one to one, we note that if the output from $G$ is $0$ or less, then it must be of the form $-\frac{n}{2}$ for an even number $n$.  Moreover, if $-\frac{n}{2} = -\frac{m}{2}$, then $n = m$. In the case we are looking at outputs from $G$ which are greater than $0$, then it must be of the form $\frac{n+1}{2}$ for an odd number $n$.  In this, too, if  $\frac{n+1}{2} = \frac{m+1}{2}$, then $n = m$.  Hence, in any case, if the outputs from $G$ are the same, then the inputs to $G$ had to be the same so  $G$ is one-to-one and $G^{-1} = \{ (0,0), (1,1), (-1,2), (2,3), (-2,4), (3,5), (-3,6), \ldots \}$  


\item  To show $P$ is a function we note that if we have the same inputs to $P$, say $2t^{5} = 2u^{5}$, then $t = u$.  Hence the corresponding outputs, $2t-1$ and $3u-1$, are equal, too. To show $P$ is one-to-one, we note that if we have the same outputs from $P$, $3t-1 = 3u-1$, then $t = u$.  Hence, the corresponding  inputs $2t^5$  and $2u^5$ are equal, too. Hence $P$ is one-to-one and $P^{-1} = \{ (3t-1, 2t^5) \, | \, \text{$t$ is a real number.} \}$

\item  To show $Q$ is a function, we note that if we have the same inputs to $Q$, say $n = m$, then the outputs from $Q$, namely $n^2$ and $m^2$ are equal. To show $Q$ is one-to-one, we note that if we get the same output from $Q$, namely $n^2 = m^2$, then $n = \pm m$.  However since $n$ and $m$ are \textit{natural} numbers, both $n$ and $m$ are positive so $n = m$. Hence $Q$ is one-to-one and $Q^{-1} = \{ (n^2, n) \, | \, \text{$n$ is a \textit{natural} number.} \}$.

\end{enumerate}

\setcounter{HW}{\value{enumi}}
\end{enumerate}


\begin{multicols}{2}
\begin{enumerate}
\setcounter{enumi}{\value{HW}}

\item $y = f^{-1}(x)$. Asymptote: $x = 0$.

\begin{mfpic}[18]{-0.5}{5.5}{-3.5}{3.5}
\axes
\tlabel[cc](5.5,-0.5){\scriptsize $x$}
\tlabel[cc](0.5,3.5){\scriptsize $y$}
\ymarks{-2, -1, 0, 1, 2}
\xmarks{ 0, 1, 2, 3,4,5}
%\tcaption{Asymptote: $x = 0$.}
\tlpointsep{4pt}
\scriptsize
\tlabel[cc](1.5, -0.5){$(1,0)$}
\tlabel[cc](1.75, 1.5){$(2,1)$}
\tlabel[cc](4, 2.5){$(4,2)$}
\axislabels {y}{{\scriptsize $-2$} -2,{\scriptsize $-1$} -1,{$1$} 1, {$2$} 2}
\axislabels {x}{{$3$} 3,{$4$} 4,{$5$} 5}
\normalsize
\penwd{1.25pt}
\arrow \reverse \arrow \parafcn{-2.5, 2.5, 0.1}{(2**t, t)}
\point[4pt]{(1,0), (2,1), (4,2)}
\end{mfpic}




\item  $y = g^{-1}(t)$. Asymptote: $y=2$.

\begin{mfpic}[9][18]{-6}{6}{-4}{3}
\axes
\tlabel[cc](6,-0.5){\scriptsize $t$}
\tlabel[cc](0.5,3){\scriptsize $y$}
\ymarks{ -3,-2,-1,1,2}
\xmarks{-5,-4,-3,-2,-1,1,2,3,4,5}
%\tcaption{Asymptote: $y=2$.}
\tlpointsep{4pt}
\scriptsize
\dashed \polyline{(-6,2), (6,2)}
\tlabel[cc](2.5, 0.5){$(2,0)$}
\tlabel[cc](1,1.25){$(0,1)$}
\tlabel[cc](2.25, -2){$(4,-2)$}
\axislabels {y}{{\scriptsize $-3$} -3,{\scriptsize $-2$} -2,{\scriptsize $-1$} -1,  {$2$} 2 }
\axislabels {x}{{$-2 \hspace{7pt}$} -2, {$-3 \hspace{7pt}$} -3,{$-4 \hspace{7pt}$} -4,  {$-5\hspace{7pt}$} -5,  {$1$} 1,  {$3$} 3, {$5$} 5, {$4$} 4}
\normalsize
\penwd{1.25pt}
\arrow \reverse \arrow \parafcn{-2.5, 2.5, 0.1}{(2*t, 2-2**t)}
\point[4pt]{(0,1), (2,0), (4,-2)}
\end{mfpic}

\setcounter{HW}{\value{enumi}}
\end{enumerate}
\end{multicols}

\begin{multicols}{2}
\begin{enumerate}
\setcounter{enumi}{\value{HW}}

\item $y = S^{-1}(t)$. Domain $[-3,3]$.

\begin{mfpic}[15]{-4}{4}{-5}{5}
\axes
\tlabel[cc](4,-0.25){\scriptsize $t$}
\tlabel[cc](0.25,4){\scriptsize $y$}
\tlabel[cc](-3,-4.5){\scriptsize $(-3,-4)$}
\tlabel[cc](0.75,-0.5){\scriptsize $(0,0)$}
\tlabel[cc](3, 4.5){\scriptsize $(3,4)$}
\ymarks{-4,-3,-2,-1,1,2,3,4}
\xmarks{-3,-2,-1,1,2,3}
\tlpointsep{5pt}
%\tcaption{Domain: $[-3,3]$.}
\scriptsize
\axislabels {y}{{$-4$} -4,{$-3$} -3,{$-2$} -2,{$-1$} -1,{$2$} 2,{$3$} 3,{$4$} 4, {$1$} 1,}
\axislabels {x}{{$-3 \hspace{7pt}$} -3,{$-2 \hspace{7pt}$} -2, {$-1 \hspace{7pt}$} -1,  {$2$} 2, {$3$} 3}
\normalsize
\penwd{1.25pt}
\parafcn{-4,4,0.1}{(3*sin(1.570796327*t/4), t)}
\point[4pt]{(-3,-4), (0,0), (3,4)}
\end{mfpic} 

\columnbreak

\item  $y = R^{-1}(s)$.  Asymptotes: $s = \pm 3$.

\begin{mfpic}[15]{-4}{4}{-5}{5}
\axes
\tlabel[cc](4,-0.25){\scriptsize $s$}
\tlabel[cc](0.25,5){\scriptsize $y$}
\gclear\tlabelrect(-1,-1.5){\scriptsize $\left(-\frac{3}{2},-\frac{1}{2} \right)$}
\tlabel[cc](-0.75,0.5){\scriptsize $(0,0)$}
\tlabel[cc](1.5,1.5){\scriptsize $\left(\frac{3}{2},\frac{1}{2} \right)$}
%\tlabel[cc](3, 3.5){\scriptsize asymptote $y=3$}
%\tlabel[cc](-2.75,-3.5){\scriptsize asymptote $y=-3$}
%\tcaption{Asymptotes: $s = \pm 3$.}
\ymarks{-4,-3,1,2,3,4}
\xmarks{-3,-2,-1,1,2,3}
\tlpointsep{5pt}
\scriptsize
%\axislabels {x}{{$-4 \hspace{7pt}$} -4,{$-3 \hspace{7pt}$} -3, {$-1 \hspace{7pt}$} -1,{$1$} 1,{$3$} 3,{$4$} 4}
%\axislabels {y}{{$-4$} -4,{$-3$} -3,{$-2$} -2, {$-1$} -1, {$1$} 1, {$2$} 2, {$3$} 3, {$4$} 4}
\normalsize
\dashed \polyline {(3, -5), (3, 5)}
\dashed \polyline {(-3,-5), (-3,5)}
\penwd{1.25pt}
\arrow \reverse \arrow \parafcn{-2.8,2.8,0.1}{( t, 0.5*(tan( 0.5236*t)))}
\point[4pt]{(0,0), (1.5,0.5), (-1.5,-0.5)}
\end{mfpic} 

\setcounter{HW}{\value{enumi}}
\end{enumerate}
\end{multicols}

\enlargethispage{1in}

\begin{enumerate}
\setcounter{enumi}{\value{HW}}


\item  

\begin{enumerate}

\item $p^{-1}(x) = \frac{450-x}{15}$.  The domain of $p^{-1}$ is the range of $p$ which is $[0,450]$

\item  $p^{-1}(105) = 23$. This means that if the price is set to $\$105$ then $23$ dOpis will be sold.

\item $\left(P\circ p^{-1}\right)(x) = -\frac{1}{15} x^2 + \frac{110}{3} x - 5000$, $0 \leq x \leq 450$.  

\smallskip

The graph of $y = \left(P\circ p^{-1}\right)(x)$ is a parabola opening downwards with vertex $\left(275, \frac{125}{3}\right) \approx (275, 41.67)$.  This means that the maximum profit is a whopping $\$41.67$ when the price per dOpi is set to $\$275$.   At this price, we can produce and sell $p^{-1}(275) = 11.\overline{6}$ dOpis.  Since we cannot sell part of a system, we need to adjust the price to sell either $11$ dOpis or $12$ dOpis. We find $p(11) = 285$ and $p(12) = 270$, which means we set the price per dOpi at either $\$285$ or $\$270$, respectively.  The profits at these prices are $\left(P\circ p^{-1}\right)(285) = 35$ and  $\left(P\circ p^{-1}\right)(270) = 40$, so it looks as if the maximum profit is $\$40$ and it is made by producing and selling $12$ dOpis a week at a price of $\$270$ per dOpi.

\end{enumerate}

\addtocounter{enumi}{1}

\item Given that $f(0) = 1$, we have $f^{-1}(1) = 0$.  Similarly $f^{-1}(5) = 1$ and $f^{-1}(-3) = -1$

\addtocounter{enumi}{9}

\item  \begin{enumerate} \addtocounter{enumii}{1} \item If $b =0$, then $m = \pm 1$.  If $b \neq 0$, then $m = -1$ and $b$ can be any real number. \end{enumerate}

\end{enumerate}

\end{document}

